\documentclass[journal]{IEEEtran}
\usepackage{cite}
\usepackage{amsmath,amssymb}
\begin{document}
\title{TriSLA: A Tri-Dimensional SLA-Aware Orchestration Architecture for O-RAN Networks with Explainable AI and Blockchain Compliance}
\author{...}
\maketitle
\begin{abstract}
Network slicing in O-RAN environments requires SLA feasibility to be assessed before resources are committed, yet many orchestration stacks implicitly assume closed-loop causality across admission, core, transport, and lifecycle subsystems. This paper presents TriSLA, a digest-frozen preventive SLA-aware runtime that evaluates feasibility at request time using semantic interpretation, PRB-governed admission scoring, multidomain telemetry snapshots, explainability metadata, detached orchestration, cluster-level reconciliation, and immutable submit-time governance. Under a single pinned container digest (sha256:ca600174…), we report evidence from controlled execution tracks spanning admission (SR, NAD), core attribution (CORE), operational contention (NCM), lifecycle persistence (LIFE), and orchestration feedback boundaries (ORCH). TriSLA demonstrates RAN-dominant hybrid admission with feasibility and resource-headroom awareness (n=450 NASP-hard+ probes), transport-informed scoring without balanced multidomain causality, runtime tri-slice score differentiation under identical network-state identifiers, and persistent NSI/reservation state with periodic reconciliation (10k+ TriSLAReservation CRDs). We further characterize non-causal boundaries: no robust tri-slice admission divergence at frozen guards, no orchestration-to-decision-engine feedback, no core-driven admission under contention campaigns, and no continuous autonomous reevaluation loop. These negative results are positioned as methodological contributions that prevent overclaim and define explicit causal limits for future closed-loop designs.
\end{abstract}
\section{Introduction}
The transition toward software-defined 5G and O-RAN network slicing shifts SLA management from static planning to runtime negotiation. Slice tenants expect service profiles—eMBB, URLLC, and mMTC—to be admitted only when radio, transport, and core resources can sustain contracted KPIs. In practice, admission controllers must therefore operate preventively: they should estimate feasibility before orchestration allocates NSI objects, reserves capacity, or triggers downstream governance actions.

Much of the existing literature emphasizes reactive assurance, simulation-only validation, or architectural diagrams that imply continuous closed-loop control without runtime proof. Orchestration platforms often describe NSI lifecycle management as if reconciliation events automatically recompute admission scores or resource pressure. Similarly, multidomain observability is frequently presented as balanced causal influence across RAN, transport, and core, even when empirical campaigns show dominance of a single domain or snapshot-based decoupling.

TriSLA addresses this gap with a reproducible Kubernetes runtime whose decision semantics, pressure formula, and guard thresholds are frozen under a single digest. The research question guiding this work is: Can SLA feasibility be assessed at request time using runtime multidomain telemetry while preserving explicit causal boundaries between admission, orchestration, reconciliation, lifecycle, and governance? We answer this question substantially—not completely—by proving preventive admission behavior and by formally characterizing paths that do not close the loop.

The contributions of this paper are sixfold. First, we demonstrate preventive SLA-aware admission via score_mode with PRB hard gates under digest-frozen execution. Second, we define and operate a PRB-governed runtime model in which resource pressure is computed from telemetry snapshots at decision time (resource_pressure_v1: 0.4·PRB + 0.3·RTT + 0.3·CPU) rather than from post-orchestration state. Third, we provide semantic persistence through SEM-CSMF SQLite stores for intents, nests, and network slices. Fourth, we implement detached orchestration in which the portal executes NASP instantiation only after decision-engine admission, without callbacks into the decision path. Fifth, we integrate immutable submit-time governance via BC-NSSMF. Sixth, we deliver a cross-track negative characterization of non-causal orchestration, core, lifecycle, and multidomain paths—an explicit scientific outcome rather than an experimental failure.

The remainder of this paper is organized as follows. Section II reviews related work on SLA admission and slicing orchestration. Section III presents the TriSLA architecture and end-to-end workflow. Section IV describes the prototype implementation and runtime scope. Section V explains the evaluation methodology and audit tracks. Section VI reports results and discusses causal boundaries. Section VII states limitations and future work. Section VIII concludes.
\section{Related Work}
Network slicing has been studied extensively from architectural, orchestration, and admission-control perspectives. Surveys on 5G slicing emphasize SDN/NFV enablement, end-to-end resource isolation, and the tension between tenant-specific QoS and shared infrastructure efficiency~\cite{Afolabi2018SlicingSurvey,OrdonezLucena2017Slicing,Foukas2017Survey,Li2021E2E}. O-RAN introduces open interfaces and programmatic RAN control, yet operational platforms still require application-level admission logic that remains auditable under real deployments~\cite{Polese2022ORAN,oran_open_smart_ran,oran-survey-2023}. SLA-aware service chaining and admission frameworks address feasibility at provisioning time~\cite{Orfanos2020SFC,Sun2022SLAAdmission,Bega2021Admission,Xu2020AIAdmission}, while AI-assisted orchestration proposals often assume multidomain feedback without publishing negative paths when feedback is absent~\cite{ahmad2023ai,sharma2022aiempowered,ai-orch-2023}.

Explainable decision support and semantic intent modeling appear in parallel lines of work~\cite{Guidotti2018XAI,xai-survey,Gruber1993,Studer1998}. Blockchain-based SLA management targets immutable contracting~\cite{Zhang2020Blockchain,Ferreira2022BlockchainSLA,Javed2024Review,bao2022blockchain,choi2023slaver}, and standards bodies define reference architectures for slice management~\cite{3GPP2020NRM,3gpp28541,ETSI2021ZSM,gsma_ng116}. Platform-oriented studies such as NASP integrate 3GPP, ETSI, and O-RAN constructs in a Network-Slice-as-a-Service workflow with detailed runtime evaluation~\cite{grings2025nasp}; TriSLA differs by foregrounding preventive admission under a digest-frozen decision plane and by reporting explicit non-causal boundaries for orchestration, core, and lifecycle subsystems.

TriSLA therefore positions itself as a reproducible runtime with PRB-governed preventive admission, detached orchestration, reconciliatory cluster state, and a documented set of limits that prior slice-management papers often omit—without claiming autonomous feedback closure across admission and orchestration.
% ... remaining sections: paste from TRISLA_CAMERA_READY_MANUSCRIPT.md
\bibliographystyle{IEEEtran}
\bibliography{references_4}
\end{document}
